\documentclass[11pt]{article}
\usepackage[margin=1in]{geometry}
\usepackage{amsmath,amssymb,amsthm}
\usepackage{booktabs}
\usepackage{hyperref}

\newtheorem{theorem}{Theorem}
\newtheorem{lemma}[theorem]{Lemma}
\newtheorem{proposition}[theorem]{Proposition}
\newtheorem{corollary}[theorem]{Corollary}
\theoremstyle{definition}
\newtheorem{definition}[theorem]{Definition}
\theoremstyle{remark}
\newtheorem{remark}[theorem]{Remark}

\title{Disproof of Conjecture 00000007709:\\
the right triangle with legs $1:2$ is a convex rep-5 tile}
\author{jilint777}
\date{2026-10-04}

\begin{document}
\sloppy
\maketitle

\begin{abstract}
Conjecture 00000007709 asserts, as its first clause, that the order $k$ of a convex
rep-tile must be a perfect square or $k=2$. This is false. The right triangle with
legs in ratio $1:2$ can be dissected into $5$ congruent triangles, each similar to it
with ratio $1/\sqrt5$, and $5$ is neither a perfect square nor $2$. We give the
dissection with integer vertex coordinates and explicit similarities
$z\mapsto(a\bar z+b)/5$ with $|a|^2=5$, and prove that it is a tiling.

A self-contained Lean~4 project (core library only) defines points, convex polygons,
similarities (via squared distances) and rep-$k$ tiles from scratch. It proves that this
triangle is a rep-5 tile and derives the negation of the clause. In Lean the plane is the
rational plane $\mathbb Q^2$; we explain why the real case follows, and the
human-readable proof below is carried out directly in $\mathbb R^2$. We also treat the
reading in which ``the order'' means the least $k\ge2$. Under that reading the
$1\times\sqrt3$ rectangle is a counterexample: its least order is $3$, and this is so
even if the pieces need only be similar to the tile, not congruent to each other.
\end{abstract}

\section{The conjecture}

The statement (English version) reads:
\begin{quote}
\emph{Definition:} The rep-tile order of a polygon is the value $k$ for which the polygon
dissects into $k$ copies similar to itself. \emph{Conjecture:} The order $k$ of a convex
rep-tile must be a perfect square or $k=2$; among convex quadrilaterals the class of
rep-tile order exactly $4$ consists of parallelograms and three exceptional families,
whose parameter spaces are two-dimensional.
\end{quote}
The Chinese version says the same thing. The conjecture is a conjunction, so refuting
one clause refutes it. We refute the first clause:
\begin{equation}\label{eq:clause}
\text{if a convex polygon is a rep-}k\text{ tile, then } k \text{ is a perfect square or } k=2.
\end{equation}
We do not discuss the second clause.

\section{Definitions and readings}\label{sec:defs}

\begin{definition}\label{def:rep}
A polygon $P\subset\mathbb R^2$ (compact, with nonempty interior) is a \emph{rep-$k$ tile}
if there are closed sets $P_1,\dots,P_k$ such that:
\begin{enumerate}
\item each $P_i=f_i(P)$ for a similarity $f_i$ of the plane;
\item the $P_i$ are pairwise congruent;
\item $P_1\cup\dots\cup P_k=P$;
\item the interiors of the $P_i$ are pairwise disjoint.
\end{enumerate}
Its \emph{orders} are the $k$ for which this holds.
\end{definition}

This is the standard notion of Golomb~\cite{golomb}, as used for example in~\cite{snover}.
As usual, ``congruent'' and ``similar'' allow mirror images. Since the areas add up,
condition~2 forces every $f_i$ to have ratio $1/\sqrt k$. Conversely, if every $f_i$ has
ratio $1/\sqrt k$, then condition~2 holds automatically, because $f_j\circ f_i^{-1}$ is
an isometry mapping $P_i$ onto $P_j$.

The conjecture's wording admits several readings. We list them and say how each is handled.
\begin{enumerate}
\item[(R1)] \emph{Every order.} ``The order $k$ of a convex rep-tile'' refers to any $k$
for which the polygon is a rep-$k$ tile. So~\eqref{eq:clause} says that every order of a
convex rep-tile is a square or $2$. This is the reading used in the conjecture itself. Its
second clause puts \emph{all} parallelograms into ``the class of rep-tile order exactly
$4$''. But the parallelogram with sides $1$ and $\sqrt2$ is also a rep-2 tile, and the
$1\times\sqrt3$ rectangle is also a rep-3 tile. So ``order $4$'' there means ``is a
rep-$4$ tile'', not ``has least order $4$''. Under (R1), Theorem~\ref{thm:main} below
refutes~\eqref{eq:clause}. This is the statement formalized in Lean.
\item[(R2)] \emph{Similar, not necessarily congruent.} The conjecture's definition says
only ``$k$ copies similar to itself''. Dropping condition~2 gives more rep-$k$ tiles, so a
counterexample under (R1) is also one under (R2): our pieces are congruent.
\item[(R3)] \emph{Least order.} ``The order'' is the least $k\ge2$ for which the polygon
is a rep-$k$ tile. ($k=1$ always works, and $1$ is a square anyway.)
Proposition~\ref{prop:rect} gives a convex polygon whose least order is $3$. This holds
with congruent pieces, and also when the pieces need only be similar to the tile.
\item[(R4)] \emph{Orientation-preserving copies only.} The pieces in
Theorem~\ref{thm:main} are mirror images of the tile, which is allowed in the standard
definition. If only orientation-preserving similarities were allowed, the rectangle of
Proposition~\ref{prop:rect} still refutes the clause. Its pieces are rotated copies, and
since a rectangle is mirror-symmetric the question does not arise.
\end{enumerate}
The clause has the escape ``$k=2$'', which does not apply to $k=5$ or $k=3$. There is no
other exceptional set.

\section{\texorpdfstring{The rep-5 tiling}{The rep-5 tiling}}\label{sec:main}

Identify $\mathbb R^2$ with $\mathbb C$. Let
\[
A=(0,0),\quad B=(10,0),\quad C=(0,5),\qquad
D=(2,4),\quad M_1=(1,2),\quad M_2=(6,2),\quad M_3=(5,0),
\]
and let $T=\triangle ABC$. This is the right triangle with right angle at $A$ and legs
$|AC|=5$ and $|AB|=10$, in ratio $1:2$. Its squared side lengths are $25,100,125$. The
point $D$ is the foot of the altitude from $A$: it lies on $BC$ ($2+2\cdot4=10$), and
$AD=(2,4)$ is orthogonal to $BC$, whose direction is $(-2,1)$. The points $M_1,M_2,M_3$
are the midpoints of $AD$, $BD$ and $AB$. Define the five triangles
\[
P_1=\triangle ADC,\quad P_2=\triangle AM_3M_1,\quad P_3=\triangle M_3BM_2,\quad
P_4=\triangle M_1M_2D,\quad P_5=\triangle M_3M_2M_1 .
\]
Geometrically, the altitude $AD$ cuts $T$ into $P_1$ (ratio $1/\sqrt5$) and
$\triangle ABD$ (ratio $2/\sqrt5$). The midlines of $\triangle ABD$ cut it into four
copies of ratio $1/\sqrt5$.

\begin{lemma}\label{lem:sim}
For $i=1,\dots,5$ let $f_i(z)=(a_i\bar z+b_i)/5$ with
\begin{center}
\begin{tabular}{cccccc}
\toprule
$i$ & 1 & 2 & 3 & 4 & 5\\
\midrule
$a_i$ & $-1-2i$ & $2-i$ & $2-i$ & $2-i$ & $-2+i$\\
$b_i$ & $10+20i$ & $5+10i$ & $30+10i$ & $10+20i$ & $25$\\
$(f_i(A),f_i(B),f_i(C))$ & $(D,A,C)$ & $(M_1,M_3,A)$ & $(M_2,B,M_3)$ & $(D,M_2,M_1)$ & $(M_3,M_1,M_2)$\\
\bottomrule
\end{tabular}
\end{center}
Then $f_i$ is a similarity with ratio $1/\sqrt5$ and $f_i(T)=P_i$. In particular the
$P_i$ are pairwise congruent and similar to $T$.
\end{lemma}

\begin{proof}
$|f_i(z)-f_i(w)|=|a_i|\,|\bar z-\bar w|/5=\sqrt5\,|z-w|/5=|z-w|/\sqrt5$, since
$|a_i|^2=5$. The vertex images in the last row are direct computations. For example,
$f_1(10)=((-1-2i)\cdot10+10+20i)/5=0=A$ and $f_1(5i)=((-1-2i)(-5i)+10+20i)/5=(-10+5i+10+20i)/5=5i=C$.
An affine bijection maps the convex hull of three points onto the convex hull of their
images, so $f_i(T)=P_i$. Two similarities of the same ratio differ by an isometry, so the
$P_i$ are pairwise congruent.
\end{proof}

A nondegenerate triangle is the intersection of the three closed half-planes that are
bounded by its edge lines and contain the opposite vertex. Its interior is the
intersection of the corresponding open half-planes. For $T$ and the pieces this gives:
\begin{center}
\begin{tabular}{ll}
\toprule
triangle & closed region (interior: all inequalities strict)\\
\midrule
$T$ & $x\ge0,\quad y\ge0,\quad x+2y\le10$\\
$P_1=ADC$ & $x\ge0,\quad y\ge2x,\quad x+2y\le10$\\
$P_2=AM_3M_1$ & $y\ge0,\quad y\le2x,\quad x+2y\le5$\\
$P_3=M_3BM_2$ & $y\ge0,\quad 2x-y\ge10,\quad x+2y\le10$\\
$P_4=M_1M_2D$ & $y\ge2,\quad y\le2x,\quad x+2y\le10$\\
$P_5=M_3M_2M_1$ & $x+2y\ge5,\quad y\le2,\quad 2x-y\le10$\\
\bottomrule
\end{tabular}
\end{center}
Each line passes through the two vertices of the corresponding edge, and the opposite
vertex satisfies the inequality strictly. For instance, for $P_5$: $x+2y=5$ passes through
$M_3,M_1$ and gives $10>5$ at $M_2$; $y=2$ passes through $M_1,M_2$ and gives $0<2$ at
$M_3$; $2x-y=10$ passes through $M_3,M_2$ and gives $0<10$ at $M_1$.

\begin{theorem}\label{thm:main}
$T$ is a rep-5 tile. Hence~\eqref{eq:clause} is false: $T$ is convex, and $5$ is neither a
perfect square nor $2$.
\end{theorem}

\begin{proof}
By Lemma~\ref{lem:sim}, the pieces $P_i=f_i(T)$ are similar to $T$ and pairwise
congruent. It remains to check conditions~3 and~4 of Definition~\ref{def:rep}.

\emph{$P_i\subseteq T$.} Each vertex of each $P_i$ is one of $A,B,C,D,M_1,M_2,M_3$, and
each of these satisfies $x\ge0$, $y\ge0$, $x+2y\le10$. Since $T$ is convex, $P_i\subseteq T$.

\emph{$T\subseteq\bigcup P_i$.} Let $(x,y)\in T$, so $x\ge0$, $y\ge0$ and $x+2y\le10$.
\begin{itemize}
\item If $y\ge2x$, then $(x,y)\in P_1$.
\item Otherwise $y\le2x$. If moreover $x+2y\le5$, then $(x,y)\in P_2$.
\item Otherwise $y\le 2x$ and $x+2y\ge5$. If $y\ge2$, then $(x,y)\in P_4$.
\item Otherwise $y\le2$ as well. If $2x-y\ge10$, then $(x,y)\in P_3$, using $y\ge0$ and $x+2y\le10$.
\item Otherwise $x+2y\ge5$, $y\le2$ and $2x-y\le10$, so $(x,y)\in P_5$.
\end{itemize}

\emph{Interiors are disjoint.} For each pair $i<j$ we give an affine function $\ell$ with
$\ell\ge0$ on $P_i$ and $\ell\le0$ on $P_j$. Then the interiors lie in $\{\ell>0\}$ and
$\{\ell<0\}$ respectively. On a triangle, an affine function is nonnegative iff it is
nonnegative at the three vertices, so this is a finite check.
\begin{center}
\begin{tabular}{llp{7.2cm}}
\toprule
pairs & $\ell$ & reason\\
\midrule
$(1,2),(1,3),(1,4),(1,5)$ & $y-2x$ & $\ell\ge0$ on $P_1$; $\ell\le0$ at every vertex of $P_2,\dots,P_5$\\
$(2,3),(2,4),(2,5)$ & $5-x-2y$ & $\ell\ge0$ on $P_2$; $\ell\le0$ at every vertex of $P_3,P_4,P_5$\\
$(3,4)$ & $2-y$ & the vertices of $P_3$ have $y\in\{0,2\}$; $P_4$ has $y\ge2$\\
$(3,5)$ & $2x-y-10$ & $\ge0$ on $P_3$ and $\le0$ on $P_5$ (table above)\\
$(4,5)$ & $y-2$ & $\ge0$ on $P_4$ and $\le0$ on $P_5$ (table above)\\
\bottomrule
\end{tabular}
\end{center}
The vertex values used are as follows. $2x-y$ equals $0$ at $A,D,M_1$, $10$ at
$M_2,M_3$, and $20$ at $B$. $x+2y$ equals $5$ at $M_1,M_3$ and $10$ at $B,D,M_2$. As a
consistency check, the doubled areas are $50$ for $T$ and $10$ for each $P_i$.
\end{proof}

\begin{remark}
In general, a triangle is a rep-$k$ tile iff $k=n^2$, $k=n^2+m^2$ or $k=3n^2$~\cite{snover}.
The present example is the case $5=1^2+2^2$. Readers who know this theorem can stop here;
the point of this submission is a complete, self-contained and machine-checked proof of one
instance. The rectangle with sides $1$ and $\sqrt k$ is a rep-$k$ tile for every $k$
(Section~\ref{sec:least}). So under reading (R1) the clause fails for every non-square
$k\ge3$.
\end{remark}

\section{\texorpdfstring{The least-order reading: the $1\times\sqrt3$ rectangle}{The least-order reading: the 1 x sqrt3 rectangle}}\label{sec:least}

\begin{proposition}\label{prop:rect}
Let $R=[0,\sqrt3]\times[0,1]$. Then:
\begin{enumerate}
\item[(a)] $R$ is a rep-3 tile, with three congruent pieces that are rotated (orientation-preserving) copies of $R$;
\item[(b)] $R$ cannot be dissected into $2$ pieces similar to $R$, whether or not the pieces are congruent.
\end{enumerate}
Hence the least $k\ge2$ for which $R$ is a rep-$k$ tile is $3$, under both the congruent and the
merely-similar notion. Since $3$ is neither a square nor $2$, the clause also fails under
reading (R3), in all its variants.
\end{proposition}

\begin{proof}
(a) For $j=1,2,3$ let $g_j(z)=iz/\sqrt3+j/\sqrt3$. This is a rotation by $90^\circ$
composed with scaling by $1/\sqrt3$ and a translation. It maps $(x,y)$ to
$\bigl((j-y)/\sqrt3,\;x/\sqrt3\bigr)$, so $g_j(R)=[(j-1)/\sqrt3,\,j/\sqrt3]\times[0,1]$. These three strips
cover $R$, and their interiors are disjoint.

(b) Suppose $R=Q_1\cup Q_2$, where $Q_i=h_i(R)$ for similarities $h_i$ of ratio $r_i$,
and the interiors are disjoint. Comparing areas gives $r_1^2+r_2^2=1$, so $0<r_i<1$. Each
$Q_i$ is a rectangle with sides $r_i\sqrt3$ and $r_i$ and diagonal $2r_i<2$.

Every corner $c$ of $R$ lies in some $Q_i$. Since $c$ is an extreme point of $R$ and
$Q_i\subseteq R$, $c$ is an extreme point of $Q_i$, that is, a vertex of $Q_i$. A single
$Q_i$ cannot contain two opposite corners of $R$: they are at distance $2$, which exceeds
the diameter $2r_i$ of $Q_i$. So each $Q_i$ contains at most two corners, and the two are
adjacent. The four corners are covered, so each $Q_i$ contains exactly two adjacent
corners, and the two pairs are complementary.

If $Q_i$ contains the two endpoints of a side $s$ of $R$, it contains $s$, by convexity.
The line through $s$ supports $R$, and hence $Q_i$. So $Q_i$ meets that line in an edge of
$Q_i$ that contains $s$ and lies in $R\cap(\text{line})=s$. Thus $s$ is an edge of $Q_i$.
A long side ($\sqrt3$) is impossible, since $r_i\sqrt3<\sqrt3$ and $r_i<1<\sqrt3$. So
$Q_1$ and $Q_2$ each have one of the two short sides of $R$, of length $1$, as an edge.
Then $r_i\sqrt3=1$ or $r_i=1$, and since $r_i<1$ this means $r_1=r_2=1/\sqrt3$. Now
$r_1^2+r_2^2=2/3\ne1$, a contradiction.
\end{proof}

\section{Lean formalization}\label{sec:lean}

The project \texttt{lean4/} (Lean 4.19.0, core library only, no Mathlib) builds everything
from scratch. All decisive objects are defined explicitly: points, convex polygons and
their interiors, similarities, rep-$k$ tiles, and the claim.

\paragraph{The plane.} A point (\texttt{Pt}) is an integer triple $(x,y,w)$ with $w>0$. It
stands for the rational point $(x/w,y/w)$, and every point of $\mathbb Q^2$ has this form.
\texttt{Pt.same} is equality of rational points ($xw'=x'w$, $yw'=y'w$). The squared
distance is $\texttt{sqNum}/\texttt{sqDen}$ with
$\texttt{sqNum}(P,Q)=(x_Pw_Q-x_Qw_P)^2+(y_Pw_Q-y_Qw_P)^2$ and
$\texttt{sqDen}(P,Q)=(w_Pw_Q)^2$.

\paragraph{Polygons.} \texttt{orient A B P} is the $3\times3$ determinant of the
homogeneous coordinates. Its sign is the sign of $(B-A)\times(P-A)$, since all $w>0$.
\texttt{ConvexPolygon vs}: the list has at least $3$ vertices, and every vertex that is not
an endpoint of an edge lies strictly to the left of that edge (counterclockwise, strictly
convex position). \texttt{InPoly} (closed region) and \texttt{InPolyInt} (interior) require
$\texttt{orient}\ge0$, respectively $>0$, for every edge. Both are invariant under
\texttt{Pt.same} (\texttt{InPoly.congr}, \texttt{InPolyInt.congr}).

\paragraph{Similarities and rep-tiles.} \texttt{IsSimilarity f n m} means
$m\cdot\texttt{sqNum}(fP,fQ)\cdot\texttt{sqDen}(P,Q)=n\cdot\texttt{sqNum}(P,Q)\cdot\texttt{sqDen}(fP,fQ)$
for all $P,Q$, that is, $|f(P)f(Q)|^2=\frac nm|PQ|^2$. Such an $f$ is automatically a
well-defined map of rational points (\texttt{IsSimilarity.respects}).
\texttt{RepTile vs k} asks for a list of $k$ maps $f_i$, each with
\texttt{IsSimilarity f 1 k} (squared ratio $1/k$, so the pieces are similar to the tile and
pairwise congruent). It further requires:
\begin{enumerate}
\item every $f_i$ maps the region into itself;
\item every point of the region is $f_i(Q)$ for some $i$ and some $Q$ in the region;
\item for $i\ne j$, no point of $f_i(\mathrm{int})$ equals a point of $f_j(\mathrm{int})$.
\end{enumerate}
For a similarity, the interior of $f_i(\text{region})$ is $f_i(\text{interior})$, so these
are conditions 3 and 4 of Definition~\ref{def:rep}. Fixing the ratio at $1/\sqrt k$ only
makes the definition stricter (Section~\ref{sec:defs}).

\paragraph{The claim.}
\texttt{Claim} $:=\forall\,\texttt{vs},\ \texttt{ConvexPolygon vs}\to\forall k,\ \texttt{RepTile vs k}\to(\exists m,\ k=m\cdot m)\lor k=2$.
This is~\eqref{eq:clause} under reading (R1), and the main theorem is
\begin{center}
\texttt{theorem conjecture\_00000007709\_false : \textasciitilde Claim}
\end{center}

\paragraph{Proof structure in Lean.}
\begin{itemize}\raggedright
\item \texttt{ConjAffine} is the map $z\mapsto(a\bar z+b)/d$ in homogeneous coordinates.
\texttt{ConjAffine.sqNum\_apply} and \texttt{ConjAffine.sqDen\_apply} are polynomial
identities showing that it scales squared distances by $|a|^2/d^2$; they give
\texttt{ConjAffine.isSimilarity}.
\item The maps \texttt{f1}--\texttt{f5} are those of Lemma~\ref{lem:sim}. \texttt{f\_similar}
shows each has squared ratio $1/5$. \texttt{f\_vertices} (by \texttt{decide}) checks the
vertex images.
\item \texttt{g1}--\texttt{g5} are the inverses $g_i(z)=a_i\bar z+c_i$. \texttt{back1}--\texttt{back5}
show that $g_i(P_i)\subseteq T$ and that $f_i(g_i(P))$ is the point $P$. With \texttt{into}
($f_i(T)\subseteq P_i$), this gives $f_i(T)=P_i$.
\item \texttt{pieces\_sub}, \texttt{cover} and \texttt{disjoint} are the three parts of the
proof of Theorem~\ref{thm:main}. They are linear-arithmetic goals in the homogeneous
coordinates, closed by \texttt{omega}.
\item \texttt{T\_rep5 : RepTile T 5} assembles these facts. \texttt{five\_not\_square} and
$5\ne2$ then give the main theorem.
\item Non-vacuity: \texttt{all\_convex} shows that $T$ and all five pieces satisfy
\texttt{ConvexPolygon}. \texttt{interiors\_nonempty} exhibits an interior point (the
centroid) of $T$ and of each piece, so the disjointness condition is not vacuous.
\end{itemize}
There is no \texttt{sorry}, no \texttt{native\_decide} and no added axiom.
\texttt{\#print axioms} reports at most \texttt{propext}, \texttt{Classical.choice} and
\texttt{Quot.sound}; \texttt{omega} uses \texttt{Classical.choice}.

\paragraph{From $\mathbb Q^2$ to $\mathbb R^2$.} The Lean statement concerns rational
points. The proof of Theorem~\ref{thm:main} above is written directly in $\mathbb R^2$, and
it is the same linear case analysis that Lean checks. The real statement also follows
formally from the rational one, as follows. The maps $f_i$ are given by the same formulas
over $\mathbb R$, and the distance identity is a polynomial identity. The pieces are closed
and their union is closed. Rational points are dense in $T$, which is convex with nonempty
interior and has rational vertices. So covering the rational points of $T$ means covering
all of $T$. If two interiors met, their intersection would be a nonempty open set, which
contains a rational point; this would contradict the rational statement. Containment
$P_i\subseteq T$ follows from convexity, because the vertices are rational points of $T$.

\paragraph{What is only in this report.} Proposition~\ref{prop:rect} (readings (R3) and
(R4)) is proved here and checked numerically in \texttt{verify.py} (exactly, in
$\mathbb Q(\sqrt3)$). It is not formalized in Lean.

\section{Reproduction}
\begin{verbatim}
cd lean4 && lake build       # Lean 4.19.0, core only
cd .. && python3 verify.py   # Python 3 standard library only
pdflatex report.tex && pdflatex report.tex
\end{verbatim}
\texttt{verify.py} uses a method different from the Lean case analysis, with exact rational
arithmetic. It checks:
\begin{itemize}
\item that $f_i(A,B,C)$ are the vertices of $P_i$ and that $|a_i|^2=5$;
\item that all pieces have squared sides $5,20,25$ (one fifth of $25,100,125$);
\item that each piece lies in $T$ (its vertices lie in $T$);
\item that each pair of pieces is separated by an edge line of one of them;
\item that the areas add up.
\end{itemize}
Containment, disjoint interiors and equal total area together imply covering: the uncovered
part of $T$ is relatively open in $T$, so if it were nonempty it would have positive area.
The script also tests $14641$ rational grid points of $T$ one by one. Finally, it verifies
the rep-3 dissection of the $1\times\sqrt3$ rectangle exactly in $\mathbb Q(\sqrt3)$.

\begin{thebibliography}{9}
\bibitem{golomb} S.~W.~Golomb, Replicating figures in the plane, \emph{Math.\ Gazette} 48
(1964), 403--412.
\bibitem{snover} S.~L.~Snover, C.~Waiveris, J.~K.~Williams, Rep-tiling for triangles,
\emph{Discrete Math.} 91 (1991), 193--200.
\end{thebibliography}

\end{document}
